\documentclass[11pt,a4paper]{article}
\usepackage[T1]{fontenc}
\usepackage[utf8]{inputenc}
\usepackage{amsmath,amssymb,graphicx,hyperref,longtable,booktabs,array}
\usepackage[margin=2.2cm]{geometry}
\providecommand{\tightlist}{\setlength{\itemsep}{0pt}\setlength{\parskip}{0pt}}
\setlength{\parindent}{0pt}\setlength{\parskip}{6pt}

\begin{document}
\hypertarget{physics-informed-lag-compensation-for-cgm-project-brief-and-prototype-findings}{%
\section{Physics-informed lag compensation for CGM: project brief and
prototype
findings}\label{physics-informed-lag-compensation-for-cgm-project-brief-and-prototype-findings}}

\hypertarget{idea}{%
\subsection{Idea}\label{idea}}

Interstitial (sensor) glucose lags plasma glucose by a diffusion delay.
Model it as first-order transport and estimate plasma glucose with a
state estimator instead of a black-box forecaster:

\[\tau \frac{dC_{isf}}{dt} = C_{bg} - C_{isf} \;\Rightarrow\; C_{bg} \approx C_{isf} + \tau \frac{dC_{isf}}{dt}\]

Worked check: CGM reads 90 mg/dL, falling 1.5 mg/dL/min, \(\tau = 10\)
min -\textgreater{} \(C_{bg} \approx 90 + 10 \times (-1.5) = 75\) mg/dL.

\hypertarget{prototype-simulation-only}{%
\subsection{Prototype (simulation
only)}\label{prototype-simulation-only}}

\begin{itemize}
\tightlist
\item
  Data: \texttt{simglucose} UVA/Padova virtual patients (30; 3 excluded
  as unphysiological -\textgreater{} 27; 5 days each), meals sized from
  each patient\textquotesingle s CR and TDI with +/-bolus error. Ground
  truth = plasma glucose; CGM = subcutaneous glucose (true
  \(\tau = 1/k_{sc}\), 6.6 to 19.9 min) + simulator Dexcom noise.
\item
  Split by patient: 14 tuning / 13 test. Estimators: raw CGM, fixed-tau
  inversion, EKF (fixed tau = 10 min), EKF with online tau (+/- 4
  fingersticks/day), EKF with true tau (oracle). EKF state = {[}G,
  dG/dt, I, ln tau, AR(1) sensor noise{]}.
\end{itemize}

\hypertarget{results-13-held-out-virtual-patients-rmse-vs-plasma-mgdl}{%
\subsection{Results (13 held-out virtual patients, RMSE vs plasma,
mg/dL)}\label{results-13-held-out-virtual-patients-rmse-vs-plasma-mgdl}}

\begin{longtable}[]{@{}llllll@{}}
\toprule\noalign{}
Sensor noise scale & Raw CGM & Fixed-tau inversion & EKF fixed tau=10 &
EKF online tau + FS & EKF true tau \\
\midrule\noalign{}
\endhead
\bottomrule\noalign{}
\endlastfoot
0 (lag only) & 8.71 & 2.66 & 2.36 & 5.05 & 0.93 \\
0.3 & 9.46 & 5.30 & 5.26 & 7.48 & 5.14 \\
1 (sd \textasciitilde12 mg/dL) & 15.08 & 15.50 & 15.56 & 14.48 &
15.56 \\
\end{longtable}

Hypoglycemia (\textless70) detection delay, noise scale 0: raw 12.4 min
late; fixed-tau EKF 1.2 min; true-tau EKF -0.3 min. At noise scale 1 the
earlier alerts cost more false alarms (raw 0.36/day; fixed-tau EKF
0.87/day).

\hypertarget{findings}{%
\subsection{Findings}\label{findings}}

\begin{enumerate}
\def\labelenumi{\arabic{enumi}.}
\tightlist
\item
  With a roughly correct tau, lag compensation removes most lag error
  (RMSE -73\% at noise 0; delay 12 -\textgreater{} \textasciitilde1
  min).
\item
  At simulator-level noise (slow, lag-1 autocorrelation 0.98, sd
  \textasciitilde12 mg/dL) no estimator beats raw CGM on pooled RMSE;
  the noise dominates.
\item
  Online tau estimation fails: median error 57-82\% vs 22\% for just
  using the 10-min prior; 4 fingersticks/day did not fix it. tau is
  weakly identifiable when glucose has no dynamic model.
\end{enumerate}

\hypertarget{caveats}{%
\subsection{Caveats}\label{caveats}}

Simulation only (lag model is exactly first-order in the simulator, so
results are optimistic); process-noise grid edge hit for some EKF
variants; some virtual children had severe simulated hypoglycemia
(excluded); no real-data (OhioT1DM) validation yet.

\hypertarget{plan}{%
\subsection{Plan}\label{plan}}

\begin{enumerate}
\def\labelenumi{\arabic{enumi}.}
\tightlist
\item
  Weeks 1-2: harden simulation (model mismatch: unmodelled extra delay,
  calibration drift, compression lows; realistic noise); tune process
  noise properly.
\item
  Weeks 3-4: make tau identifiable (add meal/insulin inputs with a
  minimal glucose-insulin model as the glucose prior; joint estimation
  with fingersticks; hierarchical population prior over tau).
\item
  Weeks 5-6: validate on OhioT1DM (needs data use agreement;
  fingersticks as reference) and compare with LSTM/linear forecasters.
\item
  Week 7: write-up: error vs noise, detection lead time vs false alarms,
  tau identifiability analysis.
\end{enumerate}
\end{document}
